\section{Sample variance and standard deviation}

For observed data \(x_1,\ldots,x_n\), define
\[
\bar x=\frac1n\sum_{i=1}^n x_i.
\]

\begin{definitionbox}
The \textbf{empirical variance} is
\[
v_n=\frac1n\sum_{i=1}^n(x_i-\bar x)^2.
\]
Its square root is the empirical standard deviation.
\end{definitionbox}

This is the variance of the empirical distribution, because every observation
has empirical mass \(1/n\).

\subsection*{A computational identity}

Expanding the squares gives
\[
v_n=\frac1n\sum_{i=1}^n x_i^2-\bar x^{\,2}.
\]
The definition explains the meaning, while this form is often convenient for
calculation.

\subsection*{The corrected sample variance}

For inference, the observations are treated before measurement as random
variables \(X_1,\ldots,X_n\). Define
\[
S^2=\frac1{n-1}\sum_{i=1}^n(X_i-\overline X)^2.
\]
After measurement, the observed value is
\[
s^2=\frac1{n-1}\sum_{i=1}^n(x_i-\bar x)^2.
\]

Assume that \(X_1,\ldots,X_n\) are independent and identically distributed with
mean \(\mu\) and variance \(\sigma^2<\infty\). The identity
\[
\sum_{i=1}^n(X_i-\overline X)^2
=
\sum_{i=1}^n(X_i-\mu)^2
-
n(\overline X-\mu)^2
\]
implies
\[
\E\left[\sum_{i=1}^n(X_i-\overline X)^2\right]
=
n\sigma^2-n\Var(\overline X).
\]
Since
\[
\Var(\overline X)=\frac{\sigma^2}{n},
\]
we obtain
\[
\E\left[\sum_{i=1}^n(X_i-\overline X)^2\right]
=(n-1)\sigma^2.
\]
Therefore
\[
\E[S^2]=\sigma^2.
\]
This is why the denominator \(n-1\) is used when the aim is to estimate the
population variance without systematic downward bias.

The random empirical variance
\[
V_n=\frac1n\sum_{i=1}^n(X_i-\overline X)^2
\]
instead satisfies
\[
\E[V_n]=\frac{n-1}{n}\sigma^2.
\]
It remains the correct variance of the empirical distribution itself.

\begin{remarkbox}
The deviations from the sample mean satisfy
\[
\sum_{i=1}^n(X_i-\overline X)=0.
\]
Thus the final deviation is determined by the preceding deviations. The factor
\(n-1\) reflects this constraint.
\end{remarkbox}

\subsection*{Two related quantities}

\begin{center}
\begin{tabular}{lll}
\toprule
Quantity & Formula & Main role \\
\midrule
Empirical variance & \(\frac1n\sum(x_i-\bar x)^2\) & describes the observed data \\
Corrected variance & \(\frac1{n-1}\sum(x_i-\bar x)^2\) & estimates population variance \\
\bottomrule
\end{tabular}
\end{center}

The theoretical variance \(\sigma^2\) belongs to a probability model. The
quantities \(v_n\) and \(s^2\) are calculated from data and need not equal
\(\sigma^2\) in a finite sample.

\begin{summarybox}
The denominator \(n\) gives the variance of the empirical distribution. The
denominator \(n-1\) produces an unbiased estimator of the population variance
under independent identical sampling with finite variance.
\end{summarybox}
